\section{Обзор литературы}
\label{sec:storage-choice}

Выбор структуры хранения для таблицы
\textit{префикс~$\rightarrow$~ASN} связан не только с задачей longest
prefix match. Для рассматриваемого микросервиса важны пять свойств:
время точечного LPM-запроса, стоимость инкрементального обновления,
расход памяти, стоимость полной выгрузки таблицы и сложность реализации в
многопоточном in-memory сервисе. Поэтому далее рассматриваются не
абстрактно лучшие LPM-алгоритмы, а семейства структур, релевантные для
сочетания чтения, обновлений и выгрузки состояния.

\subsection{PATRICIA и radix-подходы}

Классическая работа Моррисона~\cite{morrison1968patricia} вводит
PATRICIA как сжатое префиксное дерево: в нём устраняются длинные цепочки
узлов с единственным потомком. Для IP-префиксов эта идея важна потому,
что адресное пространство разрежено, а обычное побитовое дерево быстро
создаёт много служебных узлов.

PATRICIA и близкие radix-подходы естественно поддерживают префиксную
семантику: поиск может идти по пути адреса с запоминанием последнего
совпавшего префикса. Их сильная сторона -- концептуальная близость к LPM
и отсутствие необходимости перебирать все возможные длины масок. Однако
для данной работы одной только LPM-семантики недостаточно. Необходимо
также обеспечить конкурентные вставки и удаления, полную выгрузку таблицы
и простой общий контракт для двух реализаций хранилища. Поэтому
PATRICIA рассматривается как базовая точка сравнения, но не как
непосредственно выбранная структура.

\subsection{LC-trie}

LC-trie Нильссона и Карлссона~\cite{nilsson1999lctries} развивает идею
сжатого дерева и использует уровневую компрессию для уменьшения глубины
поиска. Этот подход релевантен, поскольку он показывает, как можно
снижать задержку LPM-запроса в программной реализации без перехода к
специализированному оборудованию.

Ограничение LC-trie для рассматриваемой задачи связано с обновлениями.
Структуры, оптимизированные под плотную и сжатую форму дерева, часто
требуют более сложной перестройки при изменении префиксов. В
маршрутизаторном сценарии, где основная нагрузка приходится на LPM-запросы,
это приемлемо. В микросервисе, который должен применять поток BGP-событий
по одному изменению, стоимость вставки и удаления становится критичной.
Поэтому LC-trie фиксирует ориентир по времени чтения, но не закрывает
требование к простой потоковой актуализации.

\subsection{Хеширование и поиск по длинам масок}

Хеш-таблица хорошо подходит для точного поиска по ключу: вставка,
удаление и чтение в среднем имеют малую вычислительную стоимость. Для
LPM, однако, ключом является не только адрес, но и длина маски. Один из
практических вариантов состоит в хранении нормализованных префиксов в
хеш-таблицах и проверке кандидатных масок от наиболее специфичной к
менее специфичной. Подходы такого класса обсуждаются, в частности, в
работе Вальдвогеля и соавторов~\cite{waldvogel1997scalable}.

Такой вариант релевантен для данной работы потому, что он хорошо
согласуется с инкрементальными изменениями: конкретный префикс можно
добавить или удалить без перестройки всей таблицы. Ограничение состоит в
стоимости LPM-запроса, особенно для IPv6: при прямом переборе может
проверяться до 129 длин масок. Кроме того, хеш-таблица сама по себе не
даёт упорядоченной полной выгрузки, что важно для API полного снимка и
сравнения состояний.

\subsection{Controlled prefix expansion}

Controlled prefix expansion, описанный Шринивасаном и
соавторами~\cite{srinivasan1999prefixexpansion}, ускоряет поиск за счёт
контролируемого расширения множества хранимых префиксов. Идея состоит в
том, чтобы заменить часть поисковой работы заранее подготовленным
представлением данных. Этот подход важен как пример явного обмена памяти
на время LPM-запроса.

Для forwarding-plane устройств такой компромисс может быть оправдан:
память выделяется заранее, а запросы чтения доминируют над обновлениями.
Для сервиса потоковой актуализации ситуация иная. Дополнительный расход
памяти и необходимость обновлять расширенное представление усложняют
обработку BGP-событий. Поэтому controlled prefix expansion используется
в обзоре как пример подхода, ориентированного на ускорение чтения, но не выбирается в
качестве основы внутреннего хранилища.

\subsection{Tree Bitmap}

Tree Bitmap Итертона и соавторов~\cite{eatherton2004treebitmap}
представляет префиксное дерево в виде компактных битовых карт. Подход
ориентирован на быстрый LPM-запрос и эффективное размещение данных в памяти:
битовые карты позволяют быстро определить наличие дочерних узлов и
префиксов без обхода большого числа пустых позиций.

Tree Bitmap релевантен как сильный представитель структур, специально
разработанных для маршрутного поиска. Его ограничение в контексте данной
работы состоит в том, что основная инженерная задача не сводится к
максимальному ускорению LPM-запроса. Сервис должен поддерживать
инкрементальные изменения, полную выгрузку и понятную синхронизацию
шардов. Эти требования делают специализированную структуру для чтения менее
удобной, чем изменяемый индекс общего назначения.

\subsection{Poptrie}

Poptrie Асаи и Охары~\cite{asai2015poptrie} показывает, что LPM можно
существенно ускорить на CPU общего назначения при агрессивном учёте
кэш-локальности, битовых векторов и инструкции подсчёта единичных битов.
Эта работа важна для обзора потому, что задаёт современный ориентир для
программной реализации высокоскоростного LPM-запроса.

В то же время Poptrie сильнее всего раскрывается в сценарии, где таблица
маршрутов используется преимущественно для чтения. Для рассматриваемого
микросервиса ключевой режим иной: после начальной загрузки таблица должна
обновляться потоком BGP-событий. Поэтому Poptrie рассматривается как
важная альтернатива для будущей оптимизации чтения, но не как базовая
структура текущей реализации.

\subsection{Adaptive Radix Tree}

Adaptive Radix Tree, предложенный Лейсом, Кемпером и
Нойманом~\cite{leis2013art}, изначально относится к in-memory индексам
СУБД. В ART внутренние узлы имеют адаптивный формат
\texttt{Node4}, \texttt{Node16}, \texttt{Node48} и \texttt{Node256};
структура выбирает представление в зависимости от локальной плотности
ветвления. Для данной работы эта идея важна не как готовый
маршрутизаторный алгоритм, а как изменяемый байтовый индекс с
относительно компактным хранением и поддержкой вставок, удалений и
обхода.

ART не является оптимальным решением для всех LPM-сценариев. В текущей
реализации поиск по адресу выполняется через проверку кандидатных
префиксов разной длины маски, что увеличивает стоимость \texttt{Lookup},
особенно для IPv6. Тем не менее ART хорошо соответствует инженерной
постановке работы: локальные изменения отдельных префиксов не требуют
полной перестройки таблицы, дерево можно разделить на шарды, а полный
обход остаётся реализуемым через обход всех деревьев.

\subsection{Сравнение по критериям}

Качественная сводка по рассмотренным структурам приведена в
табл.~\ref{tab:lpm-matrix}. Оценки в таблице не являются
универсальным ранжированием алгоритмов; они отражают применимость к
целевому профилю данной работы: микросервис, поток BGP-обновлений,
операции \texttt{Lookup}, \texttt{Upsert}, \texttt{Delete} и
\texttt{Dump}.

\begin{table}[ht]
\centering
\caption{Качественное сравнение структур хранения по ключевым критериям}
\label{tab:lpm-matrix}
\small
\begin{tabular}{lcccccc}
\toprule
Структура & Lookup & Update & Delete & Dump & Memory & Сложность \\
\midrule
Слайс / массив          & низк. & низк. & низк. & высок. & ср. & низк. \\
Хеш-таблицы             & ср.   & высок. & высок. & низк. & ср. & ср. \\
PATRICIA / radix        & высок. & ср. & ср. & ср. & ср. & ср. \\
LC-trie                 & высок. & низк. & низк. & ср. & ср. & высок. \\
Controlled expansion    & высок. & низк. & низк. & ср. & низк. & высок. \\
Tree Bitmap             & высок. & ср. & ср. & ср. & высок. & высок. \\
Poptrie                 & высок. & низк. & низк. & ср. & высок. & высок. \\
\textbf{ART (шард.)}    & ср. & \textbf{высок.} & \textbf{высок.} & ср. & ср. & ср. \\
\bottomrule
\end{tabular}
\end{table}

Из сравнения следует, что выбор ART не означает утверждения о его
превосходстве как LPM-алгоритма во всех условиях. Он выбран как
практический компромисс для системы, где после начальной загрузки
доминирует поток инкрементальных изменений, а операции чтения и полной
выгрузки должны оставаться работоспособными и измеримо ограниченными по
времени, памяти и числу аллокаций. Именно поэтому в дальнейших главах
рассматривается шардированная реализация ART и её экспериментальное
сравнение с прежней реализацией на \texttt{ranger}.

%==============================================================
